 % For DPF96
% APSABFRM.TEX -- The APS Meeting Electronic Abstract Submission Form.
% Version 1.1.a
% % DON'T remove the following comment; it identifies the form.
% % The American Physical Society Electronic Contributed Abstract
% Submission Form.
% ---------------------------------------------------------------------
% COMMENTS are preceded by the '%' character.
% Lines (or portions of % lines) that start with a '%' are ignored by LaTeX.
% Such lines % may be left in the file during processing.
% % Lines that are to be filled in by the author start at the left
% margin with either a '\' or a '%\'. Author-provided information
% should be entered within the curly braces. Indented lines
% beginning with '%' are examples only, and should be left alone.
% ---------------------------------------------------------------------
% The following two lines are REQUIRED. Leave them intact. Note that
% the APS office requires 11 point type.
\documentstyle[11pt,apsab]{article}
\nofiles
% ---------------------------------------------------------------------
% ADMINISTRATIVE INFORMATION
% ---------------------------------------------------------------------
% The following statement identifies the meeting for which this
% abstract is being submitted. It is a required statement and must
% contain a valid meeting ID. The valid meeting ID can be found
% in the information file for the meeting to which you are
% submitting.
% \MeetingID{DPF96}
\MeetingID{PSF15}
% ---------------------------------------------------------------------
% Submitting member information for contributed abstracts:
% These four statements are required.
\SubmittingMemberSurname{}
\SubmittingMemberGivenName{}
\SubmittingMemberEmail{}
\SubmittingMemberAffil{}
% The following line is for future use. Leave it intact.
%\SubmittingMembershipNumber{}
% ---------------------------------------------------------------------
% The following statement identifies the type of presentation.
% Insert a value from the following list:
% oral
% poster
\PresentationType{}
% ---------------------------------------------------------------------
% Use the following optional statement to indicate any desired
% special instructions, such as primary or supplemental abstract,
% special sequencing,A/V requests, etc.
\SpecialInstructions{}
% ---------------------------------------------------------------------
% Supply one sorting category for this abstract. Use one of the
% category codes listed in the information file for the meeting
% to which you are submitting, and fill in the brackets from left to
% right ,so that if there are only two levels of sorting, the first
% two brackets are filled in, and the last two are blank,
% thusly: {34}{d}{}{}.
\SortCategory{}{}{}{}
% Insert either theoretical or experimental in the next set of brackets.
\DPPType{}
% ---------------------------------------------------------------------
% SCIENTIFIC INFORMATION
% ---------------------------------------------------------------------
% The following information constitutes your printed abstract,
% subject to APS formatting rules. If you are not able to view the
% abstract with a LaTeX viewer, make certain that the title, author
% list, abstract body and footnotes do not total more than 1300
% characters, including spaces.
% DO NOT change the following line.
\begin{document}
% ---------------------------------------------------------------------
% Use the following statement to indicate the title of your abstract.
% Use mixed case and do not underline.
\Title{}
% ---------------------------------------------------------------------
% Authors and their affiliations are specified in groups by
% institution.
% Use the following three statements to indicate presenting author
% and affiliation information.
\AuthorSurname{}
\AuthorGivenName{}
\AuthorAffil{}
% If there are additional authors associated with the presenting
% author's affiliation, the 'surname' and 'given name' statements
% may be repeated prior to the 'affil' statement to save abstract
% space.
% If there are more author/affiliation combinations, create
% additional sets of \Author markup groups as many times
% as required.
% ---------------------------------------------------------------------
% The following line is REQUIRED. Leave it intact.
\begin{abstract}
% Replace this comment with your abstract. Abstract text lines
% should NOT begin with '%' characters. Use the LaTeX \footnote
% command to include footnotes in the abstract.
% The following line is REQUIRED. Leave it intact.
\end{abstract}
% ---------------------------------------------------------------------
% The following line is REQUIRED. Leave it intact.
\end{document}
% ---------------------------------------------------------------------
% Submit this abstract to The American Physical Society office via
% email to abs-submit@aps.org.
The sender will receive an %
acknowledgement of this submission electronically.
% --------------------------------------------------------------------- 